% Fragmento de entrega (sin preambulo; se compila desde main.tex).
\section*{Bitácora de decisiones -- Problema 1}

Cada decisión registra qué se decidió, qué se descartó, por qué y con qué
evidencia. Las tres primeras decisiones reproducen el texto del usuario sin
cambios; las tres últimas fueron redactadas por el asistente con evidencia
propia y aceptadas por el alumno.

\subsection*{Decisión 1 -- Tratamiento de códigos raros}

Texto del usuario, sin cambios:
\begin{quote}
``(a) Qué decidí: agrupar los códigos no documentados con la categoría
``otros'': EDUCATION 0, 5 y 6 pasan a 4, y MARRIAGE 0 pasa a 3.

(b) Alternativas que descarté: eliminar esas filas del entrenamiento, o
dejarlas como categorías propias (one-hot).

(c) Por qué: con boosting, las tres variantes dan el mismo AUC de
validación dentro del ruido: agrupar 0.7789, eliminar 0.7780 y
categorías propias 0.7788, con desviación entre particiones de
$\approx 0.005$ a $0.006$. Como el AUC no decide, elegí la más simple y la que no
pierde filas. Descarté eliminar porque quita $\approx 1.3\%$ de los datos sin
ganancia y cambiaría la población que estoy auditando; descarté las
categorías propias porque dejan columnas de 14 a 54 filas.

(d) Qué pierdo: al agrupar mezclo códigos que quizá significan cosas
distintas, y que sean equivalentes a ``otros'' es un supuesto que no
puedo verificar, porque la documentación de UCI no define esos
códigos.''
\end{quote}
Evidencia: \path{02_decisiones_datos.ipynb} (secciones 1 y 3) y
\path{resultados/significado_education.csv},
\path{significado_marriage.csv} y \path{decisiones_datos.csv}.

\subsection*{Decisión 2 -- Tratamiento de PAY\_0..PAY\_6}

Texto del usuario, sin cambios:
\begin{quote}
``(a) Qué decidí: tratar PAY\_0..PAY\_6 de forma distinta según el modelo. En
boosting las dejo como números tal cual (codificación numérica). En
regresión logística las separo en indicadores de estado ($-2$, $-1$, $0$)
más una variable de meses de atraso (valores $\ge 1$).

(b) Alternativas que descarté: en boosting, la codificación por estados más
meses de atraso; en regresión logística, la codificación numérica.

(c) Por qué: en boosting las dos codificaciones dan el mismo AUC de
validación dentro del ruido (0.7789 con numérico contra 0.7782 con
estados, con desviación entre particiones de $\approx 0.006$); los árboles
pueden cortar en $-2$, $-1$ y $0$ por su cuenta, así que usé la más simple.
En regresión logística la diferencia sí es real: el AUC pasa de 0.7267
con numérico a 0.7687 con estados ($+0.042$, con desviaciones de
$0.003$ a $0.007$), porque $-2$, $-1$ y $0$ son estados distintos y no una
escala numérica, y la logística los trata como si lo fueran.

(d) Qué pierdo: al usar codificaciones distintas, la comparación entre
modelos mezcla el efecto del modelo con el de la codificación. Además,
en boosting acepto que el árbol necesite varios cortes para separar
códigos que no tienen un orden numérico real, y no sé qué significan
el 0 y el $-2$ porque la documentación de UCI no los define.''
\end{quote}
Evidencia: \path{02_decisiones_datos.ipynb} (secciones 1, 3 y 4) y
\path{resultados/significado_pay_0.csv},
\path{significado_pay_2.csv} y \path{decisiones_datos.csv}.

\subsection*{Decisión 3 -- Incluir o excluir SEX}

Texto del usuario, sin cambios:
\begin{quote}
``(a) Qué decidí: SEX no entra como variable de entrada en el modelo
principal. Se conserva aparte, solo como variable de auditoría para
analizar los resultados por grupo.

(b) Alternativa que descarté: incluir SEX como entrada del modelo.

(c) Por qué: SEX quedó fuera porque la diferencia de AUC con y sin SEX
($0.00013$) es mucho menor que la variación entre particiones de
validación cruzada ($\approx 0.006$); no aporta valor predictivo medible, así
que no justifico usar un atributo protegido. (Con SEX: 0.77891; sin
SEX: 0.77878, boosting, 5 particiones.)

(d) Qué pierdo: la comparación se hizo solo con boosting, no con regresión
logística. Que SEX no mejore el AUC no significa que quitarla elimine
las diferencias entre grupos, porque otras variables pueden funcionar
como sustitutas; eso lo evalúo en el inciso (e) de la práctica. Además,
para analizar equidad sigo necesitando saber el grupo de cada persona,
así que el sexo no desaparece del análisis, solo del modelo.''
\end{quote}
Evidencia: \path{02_decisiones_datos.ipynb} (sección 5) y
\path{resultados/decisiones_datos.csv}.

\subsection*{Evidencia posterior (07) -- Decisión 3}

Nota agregada por la libreta \path{07_sin_atributo.ipynb} (mismo split,
mismos hiperparámetros de 03, misma calibración isotónica global, umbral
$1/6$). No forma parte del texto textual del usuario. Incluir SEX como
entrada no elimina las diferencias; las amplía: la brecha de FPR pasa de
$0.042289$ a $0.082428$ ($+0.040$, IC95\% pareado $[0.031,0.048]$, sin el
cero); la brecha de FNR pasa de $-0.023922$ a $-0.053067$ ($-0.029$,
IC95\% pareado $[-0.042,-0.017]$, sin el cero); $\Delta$AUC $=-0.0003$
($0.779162$ con SEX frente a $0.779463$ sin SEX); y $\Delta$ganancia
$=+35$ unidades de $G$ ($2019-1984$), IC95\% pareado $[-36,100]$, que sí
contiene el cero (dentro del ruido). Esto confirma la advertencia del inciso
(d) del usuario: quitar SEX del modelo no elimina las diferencias entre
grupos (otras variables predicen SEX con AUC $0.650929$ en CV y $0.649996$
en prueba), y meterla como entrada las agrava. La decisión de dejar SEX
fuera del modelo principal se mantiene.

\subsection*{Decisión 4 -- Método de calibración}

\textbf{Qué decidí:} usar calibración isotónica global, con un único
calibrador aplicado a todos los solicitantes. No se calibró por grupo.

\textbf{Qué descarté:} sin calibrar, sigmoide global y calibración por grupo.
La calibración por grupo usaría SEX para calibrar, lo que se aparta de la
decisión de dejar SEX fuera como variable de auditoría, y fragmentaría el
entrenamiento.

\textbf{Por qué:} el criterio fue menor log-loss en validación cruzada
anidada dentro del 70\% (5 particiones externas; dentro de cada una,
\texttt{CalibratedClassifierCV(cv=5)}). El isotónico obtuvo $0.426337$,
frente a $0.426558$ del sigmoide y $0.427871$ sin calibrar, y ganó también en
Brier ($0.133471$) y ECE ($0.004595$). El AUC no se usó para elegir porque la
calibración busca probabilidades correctas, no reordenar casos.

\textbf{¿Se usó la prueba para elegir?} No. El 30\% de prueba se evaluó una
sola vez después de la elección, solo como reporte (ECE isotónico en prueba
$0.012525$).

\textbf{Qué pierdo:} la isotónica es un escalón flexible que puede sobreajustar
con muestras pequeñas y producir probabilidades repetidas; el sigmoide es más
estable y suave. Además, un calibrador global garantiza
$P(Y=1\mid S=s,G=g)\approx s$ solo en promedio: no corrige por sí mismo las
diferencias de FPR/FNR entre grupos, que deben tratarse a nivel de umbral.

Evidencia: \path{04_calibracion.ipynb} y
\path{resultados/04_comparacion_cv.csv},
\path{04_metricas_test.csv}, \path{04_metodo_calibracion.csv} y
\path{04_tabla_grupos_despues.csv}.

\subsection*{Decisión 5 -- Definición de equidad}

\textbf{Qué decidí:} elegir \textbf{igual FPR} como definición de equidad.

\textbf{Qué descarté:} paridad demográfica, que iguala tasas de rechazo pero
deja diferencias de errores de $1.76$ y $1.26$ puntos porcentuales en OOF, y
paridad predictiva, que iguala la PPV pero deja diferencias de FPR de $10.96$
puntos y de FNR de $-7.51$ en OOF. Igual FNR e igual FPR y FNR no se
descartaron por malas, sino porque en esta cuadrícula eligen el mismo par que
igual FPR, así que no cambian el resultado.

\textbf{Por qué:} es la única diferencia entre grupos que el bootstrap de 05
muestra como real: el intervalo de FPR no contiene el cero
($0.042289$, IC 95\% $[0.020183,0.064465]$) y el de FNR sí
($-0.023922$, IC 95\% $[-0.064162,0.015438]$). El banco debería corregir
primero la disparidad que la evidencia respalda. Además, igual FPR ataca la
violación más visible del principio de tratar igual a quien paga: rechazar al
$36.74\%$ de los hombres que no incumplieron frente al $32.51\%$ de las
mujeres (prueba).

\textbf{Qué pierdo:} al igualar FPR se despriorizan otras definiciones. La
tolerancia ($0.01$) y la cuadrícula ($0.05$--$0.60$, paso $0.01$) son
convenciones elegidas para el análisis; otra tolerancia podría dar otro par.
La igualdad de FPR no garantiza igualdad de FNR o de PPV, y su cumplimiento
debe vigilarse en carteras nuevas.

Evidencia: \path{06_intervencion.ipynb} y
\path{resultados/06_comparativa_definiciones.csv},
\path{06_pares_umbrales_igual_fpr.csv} y
\path{06_ganancia_vs_diferencia.png}.

\subsection*{Decisión 6 -- Selección de umbrales}

\textbf{Qué decidí:} usar dos umbrales, $t_H=0.19$ y $t_M=0.18$, elegidos con
predicciones fuera de muestra (\path{cross_val_predict}, $cv=5$) sobre el
70\%. El 30\% de prueba se usó solo para reportar el resultado, sin afinar los
umbrales.

\textbf{Qué descarté:} mantener el umbral único $1/6$, que conserva una
diferencia de FPR de $4.23$ puntos porcentuales que el bootstrap de 05 señaló
como real, y elegir los umbrales sobre la prueba, lo que invalidaría la
estimación del costo y de la equidad resultante.

\textbf{Por qué:} con igual FPR y tolerancia $0.01$, el par $(0.19,0.18)$ es
el que maximiza la ganancia OOF entre los factibles ($4697$ unidades de $G$,
frente a $4745$ de la línea base). En prueba el costo no se distingue de
cero: pérdida de $7$ unidades de $G$, IC95\% $[-100,116]$ (la magnitud
puntual es $0.35\%$ de la ganancia base, $0.000778$ por solicitante, pero no
es distinguible del ruido con esta muestra). Los desempates se resolvieron de
forma determinista: menor diferencia máxima, luego menor $t_H$ y menor $t_M$.

\textbf{Qué pierdo:} usar umbrales distintos por grupo es una política
explícita de trato diferenciado por atributo protegido y requiere
justificación jurídica y de gobernanza; no basta con que sea técnicamente
barata. Además, el par se eligió en una sola cuadrícula y tolerancia, y su
estabilidad ante otras tolerancias o carteras nuevas debe vigilarse. La
intervención no elimina la información de SEX: sigue siendo necesaria para
auditar y para aplicar el umbral correcto.

Evidencia: \path{06_intervencion.ipynb} y
\path{resultados/06_umbrales_elegidos.csv},
\path{06_top_pares_igual_fpr.csv},
\path{06_costo_equidad.csv} y \path{06_bootstrap_ic95.csv}.

\subsection*{Nota sobre el recuadro (e) de la teoría}

\textbf{Qué decidí:} el recuadro ``Tú decides'' de
\path{p1_teoria.tex} (inciso (e)) propone igual FPR con
$(t_H,t_M)=(0.19,0.18)$, acepta umbrales distintos por grupo con justificación
y mantiene SEX fuera del modelo principal, en consistencia con la Decisión 3.

\textbf{Por qué:} resume los resultados de \path{06_intervencion.ipynb} y
\path{07_sin_atributo.ipynb} (ver Decisiones 5 y 6) y los conecta con la
imposibilidad del inciso (b); no introduce decisiones nuevas más allá de esas
dos.

\textbf{Qué pierdo:} el recuadro es una síntesis de redacción e
interpretación; su defensa oral se apoya en las Decisiones 5 y 6.

\textbf{Estado:} decisión aceptada y vigente.

